%
% Copyright � 2026 Peeter Joot.  All Rights Reserved.
% Licenced as described in the file LICENSE under the root directory of this GIT repository.
%
%{
\input{../latex/blogpost.tex}
\renewcommand{\basename}{cosh_complex}
%\renewcommand{\dirname}{notes/phy1520/}
\renewcommand{\dirname}{notes/ece1228-electromagnetic-theory/}
%\newcommand{\dateintitle}{}
%\newcommand{\keywords}{}

\input{../latex/peeter_prologue_print2.tex}

\usepackage{peeters_layout_exercise}
\usepackage{peeters_braket}
\usepackage{peeters_figures}
\usepackage{siunitx}
\usepackage{verbatim}
%\usepackage{macros_cal} % \LL
%\usepackage{amsthm} % proof
%\usepackage{mhchem} % \ce{}
%\usepackage{macros_bm} % \bcM
%\usepackage{macros_qed} % \qedmarker
%\usepackage{txfonts} % \ointclockwise

\beginArtNoToc

\generatetitle{XXX}
%\chapter{XXX}
%\label{chap:cosh_complex}

Show that
\begin{equation}\label{eqn:cosh_complex:20}
\Be_1 \cosh\lr{\mu + i \mu_2} = \cosh\lr{\mu - i \mu_2} \Be_1,
\end{equation}
where \( i = \Be_1 \Be_2 \).

First note that
\begin{equation}\label{eqn:cosh_complex:40}
\begin{aligned}
\Be_1 e^{i \theta }
&=
\Be_1 \lr{ \cos\theta + i \sin\theta } \\
&=
\lr{ \cos\theta - i \sin\theta } \Be_1,
\end{aligned}
\end{equation}
since \( \Be_1 i = \Be_{112} = -i \Be_1 \).

So
\begin{equation}\label{eqn:cosh_complex:60}
\begin{aligned}
\Be_1 \cosh\lr{\mu + i \mu_2}
&=
\frac{\Be_1}{2} \lr{ e^{\mu + i\mu_2} + e^{-\mu -i \mu_2} } \\
&=
\inv{2} \lr{ e^{\mu} \Be_1 e^{i\mu_2} + e^{-\mu} \Be_1 e^{-i \mu_2} } \\
&=
\inv{2} \lr{ e^{\mu} e^{-i\mu_2} \Be_1 + e^{-\mu} e^{ i \mu_2} \Be_1 } \\
&=
\cosh\lr{ \mu - i \mu_2 } \Be_1
\end{aligned}
\end{equation}

%}
%\EndArticle
\EndNoBibArticle
